\section{Background and Related Work}
\label{sec:ch4-background-related}
We first review biophilic design in architecture, then examine its treatment in HCI and Interaction Design (IxD), highlighting gaps and situating this work within these fields.

\subsection{Biophilia and Its Treatment in Architecture}
Edward Wilson popularised biophilia as humanity’s evolutionary need to connect with nature, an ``innate tendency to focus on life and lifelike processes''\cite{wilson1986biophilia}. This connection extends beyond physical contact, encompassing emotional, cognitive, and ecological awareness\cite{ulrich1984view}. A greater interaction with nature strengthens this connectedness, and those with stronger connections actively seek more nature-oriented experiences~\cite{ulrich1984view}.

Architecture has long reflected this interdependence between humans and natural environments. Traditional Japanese tea houses, for instance, embody concepts like wabi-sabi, valuing authenticity, ageing, and transience, with designs that \emph{extend into the surrounding landscape}~\cite{juniper2011wabi}. However, modern industrialisation led to building practices increasingly detached from ecological context. In response, architects like Frank Lloyd Wright championed ``organic architecture''\cite{mumford1989form}, while Ian McHarg’s \emph{Design with Nature} promoted ecological integration\cite{mcharg1969design}. Stephen Kellert later formalised biophilic design into a set of principles such as natural lighting, materiality, and biomimicry that could be embedded into human-centred design frameworks~\cite{kellert1995biophilia}. Today, biophilic design is being revisited as a response to climate change, urban stress, and declining nature connectedness. Initiatives like the Living Building Challenge\sidenote{https://living-future.org/lbc/} promote regenerative design through green roofs, vertical gardens, and water recycling. Contemporary examples include Singapore’s Gardens by the Bay\sidenote{https://www.gardensbythebay.com.sg/}, Apple Park’s ecologically integrated campus\sidenote{https://www.fosterandpartners.com/projects/apple-park}, and Maggie’s Centres~\cite{tekin2023impact}, which use natural light, greenery, and views to support healing and well-being. 


\subsection{Human-Nature Interaction in HCI}
HCI has explored the intersection of technology, people, and nature through areas such as Human-Nature Interaction and Natural Interaction~\cite{hirsch2022survey}. NatureCHI studies how people use technology in natural environments~\cite{hakkila2016naturechi}, and Outdoors HCI designs for wild or semi-wild spaces~\cite{jones2018hci, mencarini2021designing} like gardens~\cite{rodgers2019hci} and farms~\cite{liu2019designing}. Recent reviews highlight strategies to enhance human-nature relationships through technology, including engaging with urban nature, bridging indoor-outdoor experiences, and supporting personal narratives~\cite{webber2023engaging, spors2023longing, mencarini2021designing}. While these approaches emphasise interaction in natural settings, our work focuses on built environments. We explore how occupants can form intrinsic, multisensory connections to nature indoors, using biophilic design to reimagine how smart buildings support nature-connected experiences.

\subsection{Biophilia in IxD for Indoor Built Environments}
Biophilic interaction design (IxD) in indoor environments typically focuses on fostering (a) connection to nature, (b) restorative experiences, and (c) ecological awareness~\cite{rao2025smart}. Designs aiming to establish \textbf{nature-connection} often tend to keep the nature of the ``connection'' broad. Digital ``windows'' simulate outdoor landscapes to foster a sense of connection with nature~\cite{wang2015intelligent, feng2019livenature}. Soft robotic flowers evoke playful, memory-rich interactions~\cite{chooi2022symbiosis, mcdonald2015nature}, while nature-inspired wallpapers use dynamic lighting and touch-activated soundscapes to extend these sensory links~\cite{buechley2010living}.

\textbf{Restorative systems} replicate nature’s psychological and physiological benefits, especially in spaces where direct access is limited~\cite{nicoly2024restorative,masters2022virtual}. For instance, plant-like movements (i.e. responding to plant-generated electrical signals) and surfaces with cellular structure patterns were explored in confined spaces for restoration~\cite{sabinson2021plant}. \emph{Nature Jar} uses digital metaphors of growth to encourage mindfulness and restoration~\cite{wang2022nature}. Other works remind people to take breaks through nature-inspired technologies, directing their minds toward nature as a means of relaxation~\cite{jung2024leaf}. Generative art has also been used to create biophilic displays, offering evolving visuals for stress relief~\cite{xing2024exploring}. Immersive experiences, such as \emph{Virtual Nature}, leverage virtual reality (VR) to simulate restorative natural environments. 

To promote \textbf{ecological awareness}, some designs embed nature-inspired cues into daily life. For example, sensor-embedded plants nudge users into taking stairs instead of elevators, saving energy~\cite{loh2010please}. Wearable devices like \emph{FloraWear}, offer tactile, mossy surfaces, to signal nature interdependence~\cite{nam2023florawear}. Playful digital games have also explored biophilic design for promoting pro-environmental behaviour in the built environment among children~\cite{cumbo2020cci}.

\subsubsection{Sensory Modalities in Biophilic Design}
Biophilic design aims to evoke feelings of being in nature through multisensory engagement. Visual elements like projected greenery or simulated natural lighting are widely used to evoke feelings of being in nature~\cite{dias2024quantifying, yildirimmultisensory, aristizabal2021biophilic}. Auditory soundscapes mimicking birdsong or flowing water have been shown to enhance nature connectedness~\cite{dias2024quantifying, yildirimmultisensory, aristizabal2021biophilic}. In contrast, modalities like touch, smell, and taste remain under-explored despite their immersive potential. Some works combine multiple senses. Touch and visual interactions were explored in \emph{Nature-bot}, which comprised robotic flower-like displays made of different materials to encourage playful interactions~\cite{mcdonald2015nature, chooi2022symbiosis}. Similarly, \emph{LiveNature} combined tactile sensations through wool (evoking contact with sheep) with ambient displays for connection to rural environments~\cite{feng2019livenature}. Few explore sensory integration: Living Wallpaper responds to touch with birdsong and motion-triggered fading prints~\cite{buechley2010living}, while one installation blended visuals, sound, and smells to evoke river landscapes~\cite{werner2014river}.  

\subsection{Gaps and Positioning}
We see that existing biophilic research in IxD often broadly focuses on fostering connections to nature or enabling experiences of restoration and ecological awareness. Much of this work is fragmented, with limited exploration of the diversity of human experiences with nature such as sensory engagement, ecological interdependence, and cultural narratives~\cite{van2001new}.

This paper draws on architectural perspectives to surface expert imaginaries and explore new possibilities for biophilic interaction in smart environments. Rather than offering precise definitions of biophilia or prescriptive technology design guidelines, we foreground conceptual and speculative insights rooted in biophilic architecture. Our goal is to broaden the discourse on biophilic interactivity in smart buildings across HCI, architecture, and IxD, and to lay a foundation for future empirical studies and design explorations.

